%% kapa: exam=1 round=28 year=2017 period=1 form=A subject=economics subject_ko=경제학원론 seq=35 number=75
%% source: Q-Net 감정평가사 기출문제 (cst00309) artlSeq=5199603 file=extracted/1008/5199603/2017년 제28회 감정평가사 1차시험 문제지/1교시 A형.pdf
%% conversion: kordoc 4.24.51 -> tools/gen1.py (auto); review=proofread-sonnet; flags=none
\begin{question}{75}
필립스(Phillips)곡선에 관한 설명으로 옳은 것은?
\begin{choices}
  \item 필립스(A. W. Phillips)는 적응적 기대 가설을 이용하여 최초로 영국의 실업률과 인플레이션 간의 관계가 수직임을 그래프로 보였다.
  \item 1970년대 석유파동 때 미국의 단기 필립스곡선은 왼쪽으로 이동되었다.
  \item 단기 총공급곡선이 가파를수록 단기 필립스곡선은 가파른 모양을 가진다.
  \item 프리드먼(M. Friedman)과 펠프스(E. Phelps)에 따르면 실업률과 인플레이션 간에는 장기 상충(trade-off)관계가 존재한다.
  \item 자연실업률가설은 장기 필립스곡선이 우상향함을 설명한다.
\end{choices}
\end{question}
\answer{3}
